\section{関連研究}

本章では，本研究に関係する既存研究を 4 つの観点から概観する．2.1 節では較正を測る指標の定義とその起源を述べる．2.2 節では LLM に確信度を言語で表明させる手法について述べる．2.3 節では多言語環境での較正を扱った研究を述べる．2.4 節では較正を改善する手法を，プロンプトによるものと学習によるものに分けて述べる．2.5 節で本研究の位置づけを示し，2.6 節で本章をまとめる．

\subsection{較正の評価指標}

\subsubsection{較正という概念と ECE}

較正は，モデルが出力する確信度と実際の正答率の一致を指す概念であり，もとは気象予報の評価から発展した．機械学習の分野では Guo ら \cite{guo2017calibration} が広く知られる．Guo らは，分類精度の高い深層ニューラルネットワークが，層の浅い従来のモデルに比べてむしろ較正が悪化していることを示した．すなわち精度と較正は別の性質であり，精度が高いモデルが自動的に信頼できる確信度を出すわけではない．

較正の程度を測る指標として広く用いられるのが ECE（Expected Calibration Error）である．ECE は確信度の値域をいくつかの区間に分割し，各区間における平均確信度と実際の正答率の差を，区間に含まれる標本数で重みづけて平均する．値が小さいほど較正が良い．区間分割を用いる推定法は Naeini ら \cite{naeini2015binning} に基づく．

ECE は解釈しやすい一方で，区間の数と確信度の分布に値が依存する．確信度が狭い範囲に偏っていると，標本の大半が少数の区間に入るため，区間ごとの差を平均するという計算が，全体の平均どうしの差に近づく．この状況では ECE の値だけから較正の良否と正答率の高低を読み分けにくい．本研究の予備実験でもこの状況が生じた（第 6 章）．

\subsubsection{Brier Score と Murphy 分解}

Brier Score は，確信度と正誤（0 または 1）の二乗誤差の平均であり，較正と識別能力の両方を同時に評価する指標である．ECE と異なり区間分割を必要とせず，確率予測に対する厳密に適正な採点規則（strictly proper scoring rule）である．

Murphy \cite{murphy1973partition} は Brier Score を 3 つの成分に分解できることを示した．すなわち信頼度（reliability），識別度（resolution），不確実度（uncertainty）である．信頼度は，確信度を区間に分けたときの各区間における確信度と正答率のずれを表す．識別度は，区間ごとの正答率が全体の正答率からどれだけ離れているかを表し，正答する問題と誤答する問題を確信度で区別できているほど大きくなる．不確実度は，評価対象となる正誤ラベルの平均のみから決まる．

ここで注意を要するのは，不確実度が正答率のみで決まるからといって，信頼度が正答率と独立になるわけではないことである．信頼度の定義には各区間の正答率が含まれる．区間が 1 つしかない極端な場合，信頼度は平均確信度と平均正答率の差の二乗に一致する．したがって分解は，較正の水準と正答率の高低を機械的に切り分ける手段ではない．また，比較する条件ごとにモデルの回答が変われば正誤ラベルも変わるため，不確実度も条件に依存する．分解の実務上の扱いについては Ferro ら \cite{ferro2012decomposition} や Pohle \cite{pohle2020murphy} に整理がある．

本研究では，分解を較正の交絡を解消する手段としてではなく，Brier Score の内訳を把握するための補助的な情報として用いる．

\subsubsection{AUROC}

AUROC（area under the receiver operating characteristic curve）は，確信度を閾値として正答と誤答を判別したときの識別性能を表す．0.5 が無作為な判別に相当し，1.0 が完全な判別に相当する．AUROC は確信度の絶対的な水準を評価せず，順序のみを評価する．すなわち，すべての確信度を一律に 0.1 下げても AUROC は変化しない．較正が悪くとも識別はできている，という状況を捉えられるため，ECE と併用されることが多い．

\subsection{確信度の言語による表明}

\subsubsection{内部確率を用いる手法との対比}

LLM の確信度を得る素朴な方法は，モデルが割り当てる確率を用いることである．Tian ら \cite{tian2023justask} は，人間のフィードバックによる強化学習（RLHF）で調整されたモデルでは，モデル内部の確率に基づく較正が崩れていることを指摘した．また商用の API では対数確率が取得できない場合も多く，実用上の制約もある．

\subsubsection{verbalized confidence}

これに対し，モデル自身に確信度を言語で答えさせる手法が提案されている．Lin ら \cite{lin2022teaching} は，モデルが不確実性を言葉で表現するよう学習させられることを示し，Kadavath ら \cite{kadavath2022know} は，モデルが自身の回答の正しさをある程度予測できることを報告した．

Tian ら \cite{tian2023justask} は，プロンプトのみで確信度を引き出す複数の方式を比較した．数値で答えさせる方式では 0.0 から 1.0 の範囲を用いており，言語表現で答えさせる方式や，複数の候補と確率を同時に出させる方式も検討している．比較の対象となるモデル側の確率は，サンプリングによって推定した条件付き確率であり，取得したトークンの対数確率をそのまま用いたものではない．同研究は，言語で表明させた確信度がこの推定確率より良好な較正を示す場合があることを報告した．本研究が比較する 3 方式は，この研究の枠組みを踏まえたものである．

Xiong ら \cite{xiong2024llmuncertainty} は確信度の引き出し方を体系的に整理し，プロンプトの与え方によって結果が大きく変わること，そして LLM が一貫して自信過剰である傾向を示した．Yang ら \cite{yang2024verbalizedscores} も verbalized confidence の性質を調べ，プロンプトの定式化に対する感度を報告している．

\subsubsection{確信度のスケールと表現形式}

確信度をどのような尺度で答えさせるかも結果に影響する．Dai ら \cite{dai2026rescaling} は，0 から 100 の尺度で答えさせると確信度が切りのよい数値に集中することを指摘し，より粗い尺度のほうが良好な結果を示す場合があると報告している．また Mielke ら \cite{mielke2022reducing} は，対話エージェントの過信を言語表現による較正で緩和する手法を示しており，数値ではなく言語で不確実性を伝える方向の研究もある．

本研究が Ling.1S として言語表現による方式を含めるのは，実用面での意味があるためである．利用者に提示する際には「0.85」という数値より「かなり自信がある」という表現のほうが直感的に伝わる可能性がある．一方で，言語表現は段階数が限られるため確信度の解像度が落ちる恐れもある．この得失を実測で確認する必要がある．

\subsubsection{自信過剰の原因に関する知見}

Seo ら \cite{seo2025advice} は ADVICE と呼ばれる手法を提案し，確信度の推定が回答の内容に依存していないことが自信過剰の主要因であると指摘した．すなわちモデルは，自分が何と答えたかを踏まえずに確信度を述べている場合がある．この知見は，回答を得たうえで改めて確信度を尋ねる 2 段階方式（Verb.2S）を設計する動機となる．

この引用には 2 つの限定がある．第一に，ADVICE 自体はモデルの重みを更新する手法であり，本研究が対象とするプロンプトのみの手法ではない．本研究が引用するのは，自信過剰の原因に関する診断的な知見に限られる．第二に，回答を提示して尋ねるだけで確信度が回答内容に依存するようになる，あるいは較正が改善するとまでは言えない．本研究はこれを検証の対象として設定するにとどめる．

\subsection{多言語環境における較正}

多言語での較正を扱った研究として，Xue ら \cite{xue2025mlingconf} の MlingConf がある．同研究は英語，日本語，中国語，フランス語，タイ語の 5 言語を対象に，モデル内部の確率を用いる手法，p(True) を用いる手法，verbalized confidence の 3 手法を比較した．扱うタスクは，言語によらず共通の内容を問うものと，各言語圏の文化や地理に固有の内容を問うものに分かれている．また付録では，思考の連鎖を用いる条件や質問を言い換える条件など，プロンプトを変えた比較も行われている．主要な実験の対象モデルは GPT-3.5 と Llama-3.1 であり，追加の実験ではLlama-2，Vicuna，GPT-4o も扱われている．

MlingConf は日本語を対象に含む数少ない研究であり，本研究にとって最も近い先行研究である．同研究の報告によれば，言語によって較正の質は異なり，英語での較正が他言語より良好な傾向がある．たとえば GPT-3.5 を用いた TriviaQA では，英語の ECE が 16.52 であるのに対し日本語では 34.49 と報告されている．

本研究との違いは，比較している条件の組み合わせにある．MlingConf が扱うのは主として推定手法の違いであるのに対し，本研究は verbalized confidence の内部で，回答と確信度を同時に出すか別のターンに分けるか，および確信度を数値で答えさせるか言語表現で答えさせるかという組み合わせを比較する．本研究はこれを日本語版 MGSM という単一のデータセット上で行う．

\subsection{較正の改善手法}

\subsubsection{学習を伴う手法}

較正を改善する手法として，モデルの重みを更新するものがある．Li ら \cite{li2025conftuner} のConfTuner は，確信度を適切に表明するようモデルを訓練する手法である．またXiao ら \cite{xiao2025restoringcalibration} は，事後学習によって損なわれた較正を回復させる手法を報告している．これらは有効である一方，モデルの重みへの接触を要するため，商用 API経由でのみ利用できるモデルには適用できない．

\subsubsection{出力を後処理する手法}

出力された確率を事後的に補正する手法も用いられる．温度スケーリングが代表例であり，Guo ら \cite{guo2017calibration} が分類モデルに対して有効性を示した．LLM に対してもXie ら \cite{xie2024adaptivetemp} が適応的な温度調整を報告している．すべての標本に共通の単調変換を施す後処理であれば確信度の順序は保たれるため，AUROC のような順序のみを評価する指標は変化しない．一方，入力に応じて補正の強さを変える適応的な手法では，標本間の順序が保たれるとは限らない．いずれの場合も，補正のパラメータを決めるための検証データが必要になる．

\subsubsection{プロンプトによる手法}

本研究が対象とするのは，モデルの重みにも出力の後処理にも触れず，入力するプロンプトのみで較正を改善する方向である．この方向には 3 つの利点がある．第一に，重みにアクセスできない商用モデルにも適用できる．第二に，追加の学習コストがかからない．第三に，モデルが更新されても手法自体は使い続けられる．黒箱環境での較正手法については Xie ら \cite{xie2024blackboxsurvey} に整理がある．

一方で，プロンプトによる手法は 2.2.2 項で述べたとおり文面の細部に結果が左右されやすく，効果の大きさもモデルやタスクに依存する．どのような設計が有効かは，対象とする条件のもとで実測して確かめるほかない．

\subsection{本研究の位置づけ}

以上を踏まえ，本研究の位置づけを 3 点に整理する．

\textbf{第一に，日本語のタスクで引き出し方式を並べて比較する点である．} 方式の比較そのものは Tian ら \cite{tian2023justask} や Xiong ら \cite{xiong2024llmuncertainty} が英語で行っており，新しい試みではない．また MlingConf は日本語を対象に含んでいる．本研究が行うのは，claude-sonnet-4-6 と日本語版 MGSM という具体的な条件のもとで，出力のタイミングと確信度の表現形式を組み合わせた 3 条件を同一の評価手続きで比較することである．

\textbf{第二に，指標の値がどこまで較正を表しているかを検討する点である．} 既存研究の多くは ECE を主要な指標としている．しかし 2.1.1 項で述べたとおりECE の値は確信度の分布に左右され，2.1.2 項で述べたとおり Murphy 分解も較正と正答率を機械的に切り分けるわけではない．本研究は複数の指標と確信度の分布を併記し，各指標から何が言えて何が言えないかを明示する．

\textbf{第三に，較正の改善をプロンプト設計として扱う点である．} 2.4 節で述べたとおり，学習や後処理による手法は適用条件に制約がある．本研究はプロンプトの設計のみで較正を改善する方向を追求する．この方向自体は先行研究にもあるため，新規性は改善手法の具体的な設計と，その日本語タスクでの検証に置く．

なお，本研究は「日本語での較正研究が存在しない」という主張をするものではない．MlingConf が日本語を対象に含めている以上，そのような主張は正確ではない．

\subsection{本章のまとめ}

本章では，較正の評価指標，確信度の言語による表明，多言語での較正，較正の改善手法について概観した．ECE の値が確信度の分布に左右されること，およびMurphy 分解の信頼度成分も各区間の正答率を含むため，較正と正答率を機械的に分離する手段にはならないことを述べた．また，日本語を対象とした先行研究として MlingConf があり，本研究はそれとは異なる条件の組み合わせを扱うことを示した．次章では，これらを踏まえた本研究の設計方針を述べる．
